\subsection{Code-as-Policy: Robots that Write their Own Skills}
\label{sec:code}

The organising contribution of this survey is to arrange code-as-policy methods not by
application or robot embodiment, but by \emph{how much of their own competence a system
produces at run time}. Concretely, we ask five questions in sequence: does the agent (i)
write control code at all, (ii) repair that code from execution feedback within a task,
(iii) remember validated code across tasks, (iv) search a population of programs rather
than following a single trajectory of repairs, and (v) close all of these into one
open-ended loop? Each ``yes'' is a rung on the ladder shown in the \S3.1 branch of
Figure~\ref{fig:taxonomy_main}, and
the population thins sharply toward the top: of the thirteen zero-shot systems, only a
handful ever reach the combined feedback, memory, and search regime. The purpose of this
arrangement is to name the progression and to make explicit how sparsely the top rung is
populated. We use the operational definitions of Table~\ref{tab:defs} throughout: a system
exhibits feedback, memory, or search only under the conditions stated there.

\input{tables/tab_definitions}

\input{figures/fig_architectures}

\subsubsection{Zero-shot synthesis}
\label{ssec:zeroshot}
\emph{Membership condition:} a system is on this rung if and only if it emits executable control
code but closes no execution loop back to code generation, so feedback, memory, and search are
all absent. The foundational move is to treat the language model as a \emph{program synthesiser} over
a fixed library of perception and control primitives. \emph{Code-as-Policies}
\citep{codeaspolicies} showed that an LLM prompted with an API and a natural-language
instruction can emit executable Python that composes those primitives, recursively
defining undefined functions and parameterising control with arithmetic and feedback
logic. \emph{ProgPrompt} \citep{progprompt} situates the generated program in the scene
by exposing available actions and objects as importable symbols, while \emph{VoxPoser}
\citep{voxposer} has the model write code that \emph{composes 3D value maps} for a motion
planner rather than calling fixed skills, loosening the dependence on a hand-authored API.
Subsequent work broadened the input modality and task horizon; \emph{RoboCodeX}
\citep{robocodex} conditions on multimodal observations, \emph{Instruct2Act}
\citep{instruct2act} maps instructions to perception--action code, and video- and
demonstration-conditioned variants such as \emph{RoboPro} \citep{robopro},
\emph{Demo2Code} \citep{demo2code} and \emph{Text2Motion} \citep{text2motion} recover
programs from richer context. \emph{Statler} \citep{statler} maintains an explicit world
state across steps, and \emph{TidyBot} \citep{tidybot} and \emph{ChatGPT for Robotics}
\citep{chatgptrobotics} demonstrate personalisation and prompt design for real hardware.
What unites this rung, and limits it, is that the program is written \emph{once}: there is
no channel from execution back to the synthesiser, so a plan that fails at run time simply
fails.

\subsubsection{Closed-loop self-repair}
\label{ssec:selfrepair}
\emph{Membership condition:} feedback is present but memory and search are absent; the agent
revises code within a task from an execution-grounded signal yet retains nothing across tasks
and never compares a population of candidates. The second rung adds a within-task feedback loop. \emph{Inner Monologue}
\citep{innermonologue} feeds success detectors, scene descriptions and human feedback back
into the model as language, letting it replan when a step fails. \emph{DoReMi}
\citep{doremi} detects misalignments between plan and execution and triggers recovery, and
\emph{REFLECT} \citep{reflect} builds a hierarchical, multimodal summary of past
interaction that an LLM queries to explain a failure and propose a correction. More recent
systems push the monitor into the perception loop; \emph{Code-as-Monitor}
\citep{codeasmonitor} compiles spatio-temporal constraints into vision-language checks, and
\emph{AHA} \citep{aha} trains a model to reason explicitly about failure modes. The novelty
of this rung is the closed loop, but it is a \emph{short} loop: corrections are discarded
at task boundaries, so nothing accumulates.

\subsubsection{Skill-library accumulation}
\label{ssec:skilllib}
\emph{Membership condition:} memory is present; validated code is written to a cross-task store
and retrieved later, independently of whether within-task feedback is used. The third rung makes the loop persistent by writing validated code into a growing library
that later tasks retrieve and compose. Because this ``memory'' axis is a full technique
family in its own right, we catalogue its systems once, in \S\ref{sec:skill}
(\S3.4.3); \emph{Voyager} \citep{voyager} is the canonical example, curating an
ever-growing library of verified skill programs with when-to-apply guards, while related
methods distil language corrections into retrievable knowledge for future tasks~\citep{droc}. The key
distinction from \S3.1.2 is temporal scope: competence compounds \emph{across} tasks rather
than being rebuilt each time.

\subsubsection{Evolutionary program search}
\label{ssec:search}
\emph{Membership condition:} search is present; the agent maintains and selects among more than
one candidate program rather than repairing a single one. The fourth rung replaces a single chain of repairs with population-based search over
programs. Rather than iteratively debugging one candidate, the agent maintains several,
scores them by execution, and mutates the best. \emph{CaP-X} \citep{capx} provides an
interactive gym and benchmark for exactly this style of coding agent, and \emph{RoboEvolve}
\citep{roboevolve} couples a planner and a learned simulator into a co-evolutionary loop
that mines near-miss failures to stabilise search. The novelty is the shift from
\emph{repair} to \emph{search}: exploration is explicit and parallel, which trades compute
for a better chance of escaping local failures.

\subsubsection{The full self-improving loop}
\label{ssec:fullloop}
\emph{Membership condition:} all three mechanisms are present in one closed loop, so the rung
is the conjunction feedback \emph{and} memory \emph{and} search (Figure~\ref{fig:architectures}).
Only a few very recent systems satisfy it. \emph{ASPIRE} \citep{aspire} pairs a robot execution
engine that emits per-primitive multimodal traces with a skill library of validated code and a
population search over candidate programs; \emph{ENPIRE} \citep{enpire} closes an analogous loop
on physical hardware through automatic reset, parallel rollouts, and log-driven revision; and
\emph{RoboClaw} \citep{roboclaw} unifies data collection, policy learning, and execution under a
single controller with self-resetting loops. That this cell is so sparsely populated, rather
than any particular system within it, is the structural observation the survey is organised to
make.

Table~\ref{tab:compare_31} makes the funnel explicit, tabulating each system's use of
feedback (F), memory (M), and search (S).
% Table 4 — code-as-policy self-improvement matrix (Fig. 2), by rung.
{\footnotesize
\begin{longtable}{@{}p{0.36\textwidth}c ccc c@{}}
\caption{The 30 code-as-policy systems of the \S3.1 branch of Figure~\ref{fig:taxonomy_main} on the three
self-improvement mechanisms, \textbf{F}eedback, \textbf{M}emory, \textbf{S}earch, which are
determined by the rung. \yy\ present, \nn\ absent. \textbf{Self-impr.}: \yy\ full loop,
\pp\ partial, \nn\ none.}\label{tab:compare_31}\\
\toprule
\textbf{System} & \textbf{Year} & \textbf{F} & \textbf{M} & \textbf{S} & \textbf{Self-impr.}\\
\midrule
\endfirsthead
\multicolumn{6}{@{}l}{\footnotesize\emph{Table~\ref{tab:compare_31} (continued)}}\\
\toprule
\textbf{System} & \textbf{Year} & \textbf{F} & \textbf{M} & \textbf{S} & \textbf{Self-impr.}\\
\midrule
\endhead
\bottomrule
\endlastfoot
\multicolumn{6}{@{}l}{\textbf{\S3.1.1 Zero-shot synthesis} \hfill F\,\nn\ \ M\,\nn\ \ S\,\nn}\\
\textbf{Code-as-Policies}~\cite{codeaspolicies} & 2023 & \nn & \nn & \nn & \nn\\
\textbf{ProgPrompt}~\cite{progprompt} & 2023 & \nn & \nn & \nn & \nn\\
\textbf{VoxPoser}~\cite{voxposer} & 2023 & \nn & \nn & \nn & \nn\\
\textbf{Instruct2Act}~\cite{instruct2act} & 2023 & \nn & \nn & \nn & \nn\\
\textbf{ChatGPT for Robotics}~\cite{chatgptrobotics} & 2023 & \nn & \nn & \nn & \nn\\
\textbf{TidyBot}~\cite{tidybot} & 2023 & \nn & \nn & \nn & \nn\\
\textbf{RoboCodeX}~\cite{robocodex} & 2024 & \nn & \nn & \nn & \nn\\
\textbf{RoboScript}~\cite{roboscript} & 2024 & \nn & \nn & \nn & \nn\\
\textbf{Text2Motion}~\cite{text2motion} & 2023 & \nn & \nn & \nn & \nn\\
\textbf{Demo2Code}~\cite{demo2code} & 2023 & \nn & \nn & \nn & \nn\\
\textbf{Statler}~\cite{statler} & 2024 & \nn & \nn & \nn & \nn\\
\textbf{Prompt2Walk}~\cite{prompt2walk} & 2024 & \nn & \nn & \nn & \nn\\
\textbf{RoboPro}~\cite{robopro} & 2025 & \nn & \nn & \nn & \nn\\
\midrule
\multicolumn{6}{@{}l}{\textbf{\S3.1.2 Closed-loop self-repair} \hfill F\,\yy\ \ M\,\nn\ \ S\,\nn}\\
\textbf{Inner Monologue}~\cite{innermonologue} & 2022 & \yy & \nn & \nn & \nn\\
\textbf{DoReMi}~\cite{doremi} & 2024 & \yy & \nn & \nn & \nn\\
\textbf{REFLECT}~\cite{reflect} & 2023 & \yy & \nn & \nn & \nn\\
\textbf{Code-as-Monitor}~\cite{codeasmonitor} & 2025 & \yy & \nn & \nn & \nn\\
\textbf{AHA}~\cite{aha} & 2024 & \yy & \nn & \nn & \nn\\
\textbf{Introspective Planning}~\cite{introspective} & 2024 & \yy & \nn & \nn & \nn\\
\midrule
\multicolumn{6}{@{}l}{\textbf{\S3.1.3 Skill-library accumulation} \hfill F\,\nn\ \ M\,\yy\ \ S\,\nn}\\
\textbf{Voyager}~\cite{voyager} & 2023 & \nn & \yy & \nn & \pp\\
\textbf{LRLL}~\cite{lrll} & 2024 & \nn & \yy & \nn & \pp\\
\textbf{RoboCoder}~\cite{robocoder} & 2024 & \nn & \yy & \nn & \pp\\
\textbf{Uni-Skill}~\cite{uniskill} & 2026 & \nn & \yy & \nn & \pp\\
\textbf{DROC}~\cite{droc} & 2024 & \yy & \yy & \nn & \pp\\
\midrule
\multicolumn{6}{@{}l}{\textbf{\S3.1.4 Evolutionary program search} \hfill F\,\yy\ \ M\,\nn\ \ S\,\yy}\\
\textbf{CaP-X}~\cite{capx} & 2026 & \yy & \nn & \yy & \pp\\
\textbf{RoboEvolve}~\cite{roboevolve} & 2026 & \yy & \nn & \yy & \pp\\
\textbf{Code Evolution (GEAR)}~\cite{codeevolution} & 2026 & \yy & \nn & \yy & \pp\\
\midrule
\multicolumn{6}{@{}l}{\textbf{\S3.1.5 Full self-improving loop} \hfill F\,\yy\ \ M\,\yy\ \ S\,\yy}\\
\textbf{ASPIRE}~\cite{aspire} & 2026 & \yy & \yy & \yy & \yy\\
\textbf{ENPIRE}~\cite{enpire} & 2026 & \yy & \yy & \yy & \yy\\
\textbf{RoboClaw}~\cite{roboclaw} & 2026 & \yy & \yy & \yy & \yy\\
\end{longtable}
}

